\documentclass[11pt]{article}
\usepackage[T1]{fontenc}
\usepackage[utf8]{inputenc}
\usepackage{lmodern}
\usepackage{amsmath,amssymb,bm}
\usepackage{geometry}
\usepackage{hyperref}
\geometry{margin=1in}

\title{DD-Independent Schwinger Spectroscopy in \texttt{jaxQFT}}
\author{}
\date{\today}

\begin{document}
\maketitle

\section{Scope}
This note documents a DD-independent spectroscopy path for the two-dimensional
U(1) Wilson-fermion model used as a Schwinger-model proxy in
\texttt{jaxQFT}. The goal is to support
\begin{enumerate}
  \item single-pion spectroscopy in the isovector pseudoscalar channel, and
  \item two-pion spectroscopy in the maximal-isospin channel using a momentum
  basis and a generalized eigenvalue analysis.
\end{enumerate}
The domain-decomposition measurements remain separate and are \emph{not} used
here.

\section{Single-Particle Channel}
The existing inline measurement
\texttt{pion\_2pt} computes the connected pseudoscalar correlator
\begin{equation}
  C_{\pi}(t;\,p) =
  \sum_x e^{ipx}
  \langle P^a(x,t)\,P^a(0,0)^\dagger \rangle,
\end{equation}
with the local interpolating field
\begin{equation}
  P^a(x,t) = \bar\psi(x,t)\,\gamma_5\,\tau^a\,\psi(x,t).
\end{equation}
For the Schwinger-model spectroscopy workflow, the practical use is:
\begin{itemize}
  \item generate or load an ensemble of U(1) configurations,
  \item measure \texttt{pion\_2pt} for a set of spatial momenta,
  \item analyze the resulting correlators with
  \texttt{scripts/mcmc/fit\_2pt\_dispersion.py}.
\end{itemize}
This produces the single-pion dispersion relation $E_\pi(p)$ that is later used
as the free two-particle reference.

\section{Two-Particle Channel}
\subsection{Operator Basis}
The new inline measurement \texttt{pipi\_i2\_matrix} targets the maximal-isospin
two-pion channel. The operator basis is
\begin{equation}
  O_n(t) = \pi^+(p_n,t)\,\pi^+(-p_n,t),
  \qquad
  p_n = \frac{2\pi n}{L_x},
\end{equation}
where
\begin{equation}
  \pi^+(p,t) = \sum_x e^{ipx}\,\bar d(x,t)\,\gamma_5\,u(x,t).
\end{equation}
The measurement uses a user-provided list of integers
\texttt{momenta = [}$n_0,n_1,\dots$\texttt{]} to define the basis.

\subsection{Wick Contractions}
For the $\pi^+\pi^+$ channel there are no disconnected contributions. The
correlator matrix is the difference of a direct and an exchange term,
\begin{equation}
  C_{nm}(t) = D_{nm}(t) - X_{nm}(t).
\end{equation}
Using dense all-to-all propagators on a fixed source time slice $t_{\rm src}$ and
sink time $t_{\rm snk}=t_{\rm src}+t$, the implementation evaluates the spatial
sums exactly. The direct term is built from momentum-projected single-pion
blocks, while the exchange term keeps the full quark-line crossing required for
the $I=2$ channel.

\subsection{Implementation Choices}
The measurement lives in
\texttt{jaxqft/core/measurements.py} as
\texttt{TwoPionI2MatrixMeasurement}. The current implementation deliberately
keeps the scope narrow:
\begin{itemize}
  \item it is DD-independent,
  \item it is restricted to two-dimensional lattices $(L_x,L_t)$,
  \item it currently assumes the only nontrivial momentum axis is the spatial
  axis \texttt{mom\_axis = 0},
  \item it uses the dense inverse path only,
  \item it averages over a configurable list of source time slices
  \texttt{source\_times}.
\end{itemize}
This is the correct tradeoff for a spectroscopy code path whose immediate target
is small-volume quenched studies and offline analysis.

\section{Offline GEVP Analysis}
The new analysis script
\texttt{scripts/mcmc/analyze\_pipi\_i2\_gevp.py} reads checkpoint
\texttt{state.inline\_records} produced by \texttt{mcmc.py}. Its workflow is:
\begin{enumerate}
  \item extract the correlator-matrix samples for one channel
  (\texttt{full}, \texttt{direct}, or \texttt{exchange}),
  \item optionally hermitize $C(t)$ sample by sample,
  \item estimate an integrated autocorrelation time from
  $\mathrm{Re}\,C_{00}(t_{\rm ref})$ and choose a jackknife block size
  $b \simeq 2\tau_{\mathrm{int}}$ unless the user overrides it,
  \item build blocked jackknife samples,
  \item solve the GEVP
  \begin{equation}
    C(t)\,v_n(t,t_0) = \lambda_n(t,t_0)\,C(t_0)\,v_n(t,t_0),
  \end{equation}
  \item construct effective energies from the principal correlators,
  \item estimate level energies from a user-selected fit window.
\end{enumerate}
The script writes both a PNG summary and a JSON file containing the numerical
results.

\section{Configuration Interface}
The measurement can be configured directly in the \texttt{mcmc.py} TOML file:
\begin{verbatim}
[[measurements.inline]]
type = "pipi_i2_matrix"
name = "pipi_i2_matrix"
every = 20
momenta = [0, 1, 2, 3]
mom_axis = 0
source_times = [0]
dense_max_dof = 4096
include_direct = true
include_exchange = true
include_full = true
\end{verbatim}
This path is independent of the DD measurement types.

\section{Smoke-Test Commands}
The code path was smoke-tested locally with the JAX environment in
\texttt{/opt/python/jax}. The tests exercised:
\begin{itemize}
  \item direct construction of \texttt{TwoPionI2MatrixMeasurement},
  \item the \texttt{build\_inline\_measurements} $\rightarrow$
  \texttt{run\_inline\_measurements} path used by \texttt{mcmc.py},
  \item offline GEVP analysis on a tiny synthetic checkpoint assembled from real
  U(1) Wilson measurements on a $4\times 8$ lattice.
\end{itemize}

\section{Current Limitations}
The present implementation is intentionally the first useful version rather than
the final spectroscopy system:
\begin{itemize}
  \item the two-pion matrix is dense-inverse only,
  \item the current channel is maximal-isospin $I=2$,
  \item the GEVP analysis uses blocked jackknife errors and an explicit fit
  window rather than a fully automatic multi-exponential model selection,
  \item free-vs-interacting level comparisons should use the measured
  single-pion dispersion relation rather than continuum kinematics alone.
\end{itemize}
These limitations are acceptable for the intended near-term use:
quenched pilot studies, operator-basis debugging, and student-facing offline
analysis exercises.

\end{document}
